\exer{Géométrie}
\setcounter{exo}{0}
%CCS PSI 2008


\begin{obj}
Respecter les normes de sécurité : l'effort de pincement doit être inférieur à 150~N.
\end{obj}

Le système de verrouillage doit maintenir la porte fermée sous l'action des passagers et des
actions dues aux différences de pressions induites par le système de climatisation.
Un calcul préalable a établi que ces actions peuvent être modélisées par une force du basculateur sur la bielle. Son point d'application est E lorsque la porte est verrouillée. Le modèle du dispositif retenu est donné sur la figure suivante.

\begin{center}
\includegraphics[width=\linewidth]{984_01}%
\end{center}

On considere le repere ayant pour origine I et pour base (x1, y1, z1). La direction x1 est colineaire au segment IE dans toutes
les configurations de la bielle de verrouillage FE et du bras IF lie au stator
du motoreducteur. Les notations retenues sont celles definies sur la figure 4. Les
liaisons entre la bielle de verrouillage et le basculeur et entre le stator et la poutre
de fermeture sont respectivement des liaisons pivot selon les axes (E, z1) et (I, z1).
La direction x2 est liee au stator et son angle par rapport a x1 vaut theta(t).
La direction x3 est liee a la bielle de verrouillage et colineaire au segment FE.
Pendant la phase de verrouillage, theta(t) varie quand le stator du motoreducteur
tourne par rapport a la poutre de fermeture. La vitesse angulaire est notee theta point,
avec IF = ell1 x2 et FE = ell2 x3 ; par construction ell1 est different de ell2.

\question{Exprimer x(t) en fonction de ell1, ell2 et theta(t), sous la forme x(t) = lambda1(t) ell1 + lambda2(t) ell2.}
\ifprof
\begin{corrige}
\end{corrige}
\else
\fi

\question{Exprimer la vitesse $\dot{x}$ du point $E$ appartenant au basculeur dans son mouvement
par rapport à $\rep{1}$ en fonction de $\ell_1$, $\ell_2$, $\theta(t)$ et $\dot{\theta}(t)$ sous la forme $A(\theta)\dot{x}=B(\theta)\dot{\theta}$.}
\ifprof
\begin{corrige}
\end{corrige}
\else
\fi

\question{Montrer qu'il existe une valeur de pour laquelle le basculeur ne pourra fournir aucune puissance au système. Expliquer pourquoi cette position est instable.}
\ifprof
\begin{corrige}
\end{corrige}
\else
\fi

\question{Un tel mécanisme est appelé mécanisme à genouillère. Existe-t-il dans le
système de fermeture d'autres cas d'utilisation de ce principe ? Dans l'affirmative,
indiquer pour chacun, les trois points considérés et le mouvement qui est
bloqué.}
\ifprof
\begin{corrige}
\end{corrige}
\else
\fi

\question{Expliquer quelle est la fonction des butées 1 et 2 (figures suivantes) et préciser
en la justifiant, leur position de réglage ?}
\ifprof
\begin{corrige}
\end{corrige}
\else
\fi

\question{Expliquer la pertinence de cette solution à empêcher l'ouverture des vantaux
quelle que soit la poussée des passagers. Conclure quant à la satisfaction de l'exigence.}
\ifprof
\begin{corrige}
\end{corrige}
\else
\fi


\begin{center}
\includegraphics[width=\linewidth]{984_02}%

 \textit{Étape de coulissement}
\end{center}

\begin{center}
\includegraphics[width=\linewidth]{984_03}%

 \textit{Étape verrouillée}
\end{center}


\begin{enumerate}
\item $x(t)=\ell_1 \cos\theta +\ell_2 \sqrt{1-r^2\sin^2\theta}$.
\item $A=\sqrt{1-r^2\sin^2\theta}$, $B=-\ell_1\sin\theta\sqrt{1-r^2\sin^2\theta}-\ell_2r^2\sin\theta\cos\theta$.
\item $\theta=0\degres$.
\item ...
\item ...
\end{enumerate}

% BEGIN AUTO CORRIGE INCLUDE
\InclureCorrige{corrige.tex}
% END AUTO CORRIGE INCLUDE
